\section{为什么这是“最后访谈”}

这期节目的特殊性在于，录制时收购尚未发生，发布时结局已经出现。读者因此能看到一个尚未被收购叙事改写的 Manus：它还在讨论自己的增长、风险、复杂性、Agent 市场和未来十年，而不是事后为交易寻找解释。

\subsection{事后读法和当时读法}

当时读，Peak 的焦点是产品和公司：如何增长，如何避免复杂，如何面对大厂，如何让 Agent 进入更多人群。事后读，Meta 收购让这些内容多了一层意义：一个还在强调简单和特色的 Agent 公司，被放进更大平台之后，能否保持这种特色？这也是本期最值得持续观察的问题。

\begin{knowledgebox}{最后访谈的价值}
最后访谈不是因为它预言了收购，而是因为它记录了收购前的原始产品判断。它让我们能比较：团队进入大平台后，哪些判断被保留，哪些会被组织和战略改写。
\end{knowledgebox}

\subsection{Manus 与 Meta 的潜在张力}

Meta 有模型、社交图谱、基础设施和全球产品经验；Manus 有更清晰的 Agent 产品心智和早期用户反馈。两者结合的可能性很大，但张力也明显：大平台倾向扩展功能、整合生态、服务更多入口；Manus 最害怕的正是失去特色和变复杂。

\begin{warningbox}{收购后的核心风险}
被大平台收购不自动等于产品变强。真正的挑战是：资源变多后，Manus 能否仍然克制，仍然让用户清楚知道它解决什么问题。
\end{warningbox}

\subsection{本章小结}

这期作为最后访谈，价值在于保留了 Manus 被大平台吸收前的自我描述。后续评估收购成败，应该回到这些原始判断。

